\section{Limitations}
\label{sec:limitations}

The current paper should be read as a synthetic benchmark and prototype study.
It does not use real FIR records, CCTNS logs, ICJS logs, police access logs, or
sensitive operational data. The policy oracle is a declared benchmark policy,
not an official police policy. Therefore, the results show whether the proposed
architecture is technically measurable, not whether it is ready for direct
deployment.

The blockchain layer is also limited. It is a local permissioned audit
simulation, not a live Hyperledger Fabric deployment. The metadata-exposure
result is a schema-level proxy, not a formal privacy proof. The audit layer can
support tamper evidence and reconstruction, but it does not prove that a policy
decision was correct.

NS-PI has a narrow role. It is not the best overall tamper detector and it does
not replace cryptographic audit, ABAC/PBAC enforcement, or trusted
raw-attribute policy checking. It is useful as a complementary log-only drift
signal for validly re-signed compromised-signer behavior. It still misses
low-rate and small targeted attacks in the current benchmark, especially
2\%--5\% global corruption and 10\% targeted station/district corruption.

The XAI layer is strong as a structured trace, but it is not perfect as natural
language. The benchmark records decision labels, reason codes, decisive
attributes, policy versions, counterfactual information, and explanation
hashes. However, the rendered explanation text does not always mention every
decisive attribute. This should be improved before the paper makes stronger
claims about explanation quality.

Any real-world pilot would require official policy validation, legal and
privacy review, secure integration design, institutional approval, and testing
with appropriate governance. This paper does not claim those outcomes.
